
% Much recent progress in the numerical simulation of mechanical models can be
% related to a trend known as the ``geometrization of mechanics''.
% Geometry has always played a central role in the development of
% mechanics (Euler, Lagrange, Cosserat, etc.),
% but with 
% the work of authors like Synge, Abraham and Arnold, these geometric 
% developments took a new direction as 
% % With every confirmation of the predictive power in Einstein's theory of relativity, 
% % 
% % Replace following with reference to Einstein's relativity and differential
% % geometry.
% the work of \cite{ricci1901methodes} began to be used to
% characterize the geometric structure of non-physical aspects of
% problems, like spaces of possible solutions.
% This newly found geometric structure in the abstract mathematics would
% eventually prove instrumental in the development of computational
% methods. 
% The work of Simo \citep{simo1986threedimensional,simo1990stress} marked a significant advance in this
% direction, as it directly drew from this abstract geometric structure to
% inform an enlightened computational procedure.

%
\section{Mechanical Governing Equations}\label{sec:apparatus}
Developments are centered around static equilibrium problems %that are
governed by nonlinear differential equations of the form:
\begin{equation}
\operatorname{B}_\chi \bm{\sigma} - \bm{f} = \dot{\bm{\mu}}
\label{eq:strong}
\end{equation}
where \(\operatorname{B}_\chi\) is a nonlinear
differential operator and the generalized stress \(\bm{\sigma}\), momentum rate \(\dot{\bm{\mu}}\), and body force \(\bm{f}\) are defined on a body \(\mathcal{B}\) placed in space by a
map \(\chi\). 
In general \(\operatorname{B}\) depends nonlinearly on \(\chi\),
but acts linearly on \(\bm{\sigma}\). Problems with the form of
\eqref{eq:strong} describe the equilibrium of continua, rods, and shells, whose
configurations \(\chi\) may comprise a \emph{nonlinear} space
\(\mathscr{C}\).

This chapter presents the mathematical foundation for later developments. In
\cref{box:mechanical-bvp} the mathematical structure which underpins
\eqref{eq:strong} is summarized, and in \cref{sec:derivatives} attention
is brought to some important subtleties that arise with derivatives on manifolds. 
This material is adapted from treatments by \cite{marsden1994mathematical,zienkiewicz2014finite,ogden1997nonlinear} and \cite{shilov1996elementary}.

\begin{ambox}[The Mechanical BVP]{box:mechanical-bvp}
{
\begin{description}
\item[Body]
A body \(\mathcal{B}\) is a set of points \(\bm{\xi}\), and a placement
\(\mathcal{B}_\chi\) of \(\mathcal{B}\) is a Riemannian manifold
\(\mathcal{B}_\chi\) with boundary \(\partial \mathcal{B}\) and tangent
bundle \(T\mathcal{B}_\chi\) embedded in an ambient Euclidean space
\(\mathcal{E}\). The map
\(\chi\colon \mathcal{B}\rightarrow \mathcal{B}_\chi\) associates the
state of a material point \(\bm{\xi}\) to the placement
\(\mathcal{B}_\chi\). A reference configuration
\(\chi_0\colon \mathcal{B}\rightarrow \mathcal{B}_{\chi_0}\) is often
implied and the pull-back by the composition \(\chi\circ\chi_0^{-1}\) is
denoted by \(\bm{A}\).
% \tightlist
\item[Metric]
The set \(\mathscr{C}\) of configurations \(\chi\) forms a smooth
 manifold. A configuration \(\chi\) is identified by a couple
\((\bm{x}, \bm{\Lambda})\) where \(\bm{x}\in\mathscr{C}_x\) is a
vectorial configuration variable, and
\(\bm{\Lambda}\in\mathscr{C}_{\scriptscriptstyle{\Lambda}}\) a
non-vectorial variable. These enter independently in a Riemannian metric
on the tangent space \(T_\chi\mathscr{C}\) to \(\mathscr{C}\) at
\(\chi\), which takes the assumed form: 
\begin{equation}
  \left< \cdot , \, \cdot\right>_{\chi} = \left< \cdot , \, \cdot\right>_{x} + \left< \cdot , \, \cdot\right>_{\scriptscriptstyle{\Lambda}} .
  \label{eq:conf-pair}
\end{equation}
%
\item[Stress]
The spatial stress \(\bm{\sigma}\) in \eqref{eq:strong} is related to a
material stress \(\bm{s}\) which is defined with respect to inner products 
in \(\mathscr{C}\) at the current \(\chi\) and reference
\(\chi_0\) configurations:
\begin{equation}
  \left<\bm{c},\, \bm{\sigma}\right>_{\chi}
  =\left< \bm{A}\bm{c},\, \bm{s}\right>_{\chi_0}
  = \left<\bm{c},\, \bm{A}_*\bm{s}\right>_{\chi}
  =\left< \bm{A}\bm{c},\, \bm{s}\right>_{\chi_0}\qquad\text{ for any }\bm{c},
  \label{eq:def-Astar}
\end{equation}
and where the transformation \(\bm{A}_*\) has been defined.
% with \(\bm{A}_*\) denoting the Piola transformation. 
Similarly, a dual
\(\operatorname{B}^*\) to the equilibrium operator \(\operatorname{B}\) in
\eqref{eq:strong} is defined with respect to \eqref{eq:conf-pair} for pure displacement boundary conditions:
\begin{equation}
  \left<\bm{\eta},\, \operatorname{B}_\chi\bm{\sigma}\right>_\chi = \left<\operatorname{B}_\chi^*\bm{\eta},\, \bm{\sigma}\right>_\chi
  \qquad \text{ for any }\bm{\eta}
  \label{eq:def-adjoint}
\end{equation}
%
\item[Strain]
The material strain \(\bm{e}\) is related to 
\(\operatorname{B}^*_\chi\) through the differential relation: 
\begin{equation}
  D\bm{e}\cdot\bm{u} \equiv \bm{A}\operatorname{B}^*_\chi \bm{u}.
  \label{eq:vstrain}
\end{equation}
%
\item[Constitution]
The constitutive response is defined between the material stress
\(\bm{s}\) and the material strain \(\bm{e}\) with material \(\bm{C}_s\) and
spatial \(\bm{C}_\sigma\) moduli defined by:
\begin{equation}
  D \bm{s}\cdot \bm{u} = \bm{C}_s \, D \bm{e}\cdot {\bm{u}}
  \qquad\text{ and }\qquad
  \bm{C}_\sigma \triangleq \bm{A}_*\bm{C}_s\bm{A}.
  \label{eq:def-moduli}
\end{equation}
\end{description}
}
\end{ambox}

The structure outlined in Box~\ref{box:mechanical-bvp} 
implies a weak statement
of \eqref{eq:strong} that is centered around the Galerkin \emph{residual functional}:
%
\begin{equation}
\mathcal{G}\left[\bm{\eta}\right] \triangleq \left< \bm{\eta}, \, \mathbb{B}_\chi \bm{\sigma}\right>_\chi - \left< \bm{\eta}, \, \bm{f}\right>_\chi.
\label{eq:galerkin}
\end{equation}
This informs our interpretation of the \emph{principle of virtual work}: forces are by definition resistance to variations of the configuration, and thus the principle of virtual work mandates that these belong to a vector space that is dual to the tangent space of the configuration manifold.



% \begin{exmpl}[Elastostatics.]{ex:elastostatics}
% Consider the elastostatic problem:
% $$
% \operatorname{div}\bm{\sigma} + \bm{b} = \bm{0}
% $$
% and the tensor products:
% $$
% \begin{aligned}
% & (\mathbf{A} \otimes \mathbf{B})[ \mathbf{S}]=(\mathbf{A} \otimes \mathbf{B}): \mathbf{S} = \mathbf{A} (\mathbf{B}\cdot \mathbf{S}) \\
% & (\mathbf{A} \oslash \mathbf{B})[ \mathbf{S}]=(\mathbf{A} \oslash \mathbf{B}): \mathbf{S}=\mathbf{A} (\mathbf{S} \circ \mathbf{B}^{\mathrm{t}})=\mathbf{A} \mathbf{S} \mathbf{B}^{\mathrm{t}} \\
% %& (\mathbf{A} ? \mathbf{B}): \mathbf{S}=\mathbf{A} \cdot \mathbf{S}^{\mathrm{t}} \cdot \mathbf{B}^{\mathrm{t}} \\
% %& (\mathbf{A} \odot \mathbf{B}): \mathbf{S}=\frac{1}{2}(\mathbf{A} \oslash \mathbf{B}+\mathbf{A} \ominus \mathbf{B}): \mathbf{S}=\frac{1}{2}\left(\mathbf{A} \mathbf{S} \mathbf{B}^{\mathrm{t}}+\mathbf{A} \mathbf{S}^{\mathrm{t}}  \mathbf{B}^{\mathrm{t}}\right)
% \end{aligned}
% $$
% \begin{description}
%     \item[Body] 
%     For a continuum body with reference configuration 
%     $\mathcal{B}_0 \subset \mathbb{R}^{n_{\mathrm{dim}}}$ 
%     and prescribed deformations on
%     the part of the boundary $\Gamma_{\chi} \subset \partial \Omega$, define the configuration space as the set
%     $$
%     \mathscr{C}=\left\{\chi \in W^{1, p}(\Omega)^{n_{\mathrm{dim}}}
%     \quad|\quad
%     \det[D \chi]>0 \text{ in }\Omega\text{ and }\left.\chi\right|_{\Gamma_{\chi}}=\bar{\chi}\right\}
%     $$
%     $$
%     \mathcal{V}_{\chi}=\left\{\bm{\eta}: \bm{\chi}(\mathscr{B}) \rightarrow \mathbb{R}^{n_{\mathrm{dim}}}: \bm{\eta}(\chi (\bm{\xi}))
%     =\bm{0}\text{ for }\bm{\xi} \in \Gamma_{\chi}\right\}
%     $$
%     $$
%     \mathcal{V}_0 =
%     \left\{\bm{\eta}_0 \quad | \quad 
%     W^{1, p}(\Omega)^{n_{\mathrm{dim}}}: \bm{\eta}_0(\bm{\xi})=\bm{0} \text{ for } \bm{\xi} \in \Gamma_{\chi}\right\}
%     $$
    
%     \item[Stress] 
%     $\bm{A}_{*} = J^{-1}\bm{F}\oslash\bm{F}$
%     $$
%     \mathbb{B} = -\operatorname{div}
%     \qquad\text{ and }\qquad
%     \mathbb{B}^* = \operatorname{grad}
%     $$
% \end{description}
% \end{exmpl}


\iffalse
\hypertarget{sec:statement}{%
\subsection{Problem Statement}\label{sec:statement}}

\eqref{eq:strong} is investigated in the context of the following
mathematical structure:

\begin{description}
% \tightlist
\item[A0]
The body \(\mathcal{B}\) is a set of points \(\bm{\xi}\), and a placement
\(\mathcal{B}_\chi\) of \(\mathcal{B}\) is a Riemannian manifold
\(\mathcal{B}_\chi\) with boundary \(\partial \mathcal{B}\) and tangent
bundle \(T\mathcal{B}_\chi\) embedded in an ambient Euclidean space
\(\mathcal{E}\). The map
\(\chi\colon \mathcal{B}\rightarrow \mathcal{B}_\chi\) associates the
state of a material point \(\bm{\xi}\) to the placement
\(\mathcal{B}_\chi\). A reference configuration
\(\chi_0\colon \mathcal{B}\rightarrow \mathcal{B}_{\chi_0}\) is often
implied and the pull-back by the composition \(\chi\circ\chi_0^{-1}\) is
denoted by \(\bm{A}\).
% \tightlist
\item[A1]
The set \(\mathscr{C}\) of configurations \(\chi\) forms a smooth
 manifold. A configuration \(\chi\) is identified by a couple
\((\bm{x}, \bm{\Lambda})\) where \(\bm{x}\in\mathscr{C}_x\) is a
vectorial configuration variable, and
\(\bm{\Lambda}\in\mathscr{C}_{\scriptscriptstyle{\Lambda}}\) a
non-vectorial variable. 
These enter independently in a Riemannian metric
on the tangent space \(T_\chi\mathscr{C}\) to \(\mathscr{C}\) at
\(\chi\) which takes the assumed form: 
\begin{equation}
  \left< \cdot , \, \cdot\right>_{\chi} = \left< \cdot , \, \cdot\right>_{x} + \left< \cdot , \, \cdot\right>_{\scriptscriptstyle{\Lambda}} .
  \label{eq:conf-pair}
\end{equation}
%
\item[A2]
The spatial stress \(\bm{\sigma}\) in \eqref{eq:strong} is related to a
material stress \(\bm{s}\) which is defined with respect to inner products 
in \(\mathscr{C}\) at the current \(\chi\) and reference
\(\chi_0\) configurations:
\begin{equation}
  \left<\bm{c},\, \bm{\sigma}\right>_{\chi}
  =\left< \bm{A}\bm{c},\, \bm{s}\right>_{\chi_0}
  = \left<\bm{c},\, \bm{A}_*\bm{s}\right>_{\chi}
  %=\left< \bm{A}\bm{c},\, \bm{s}\right>_{\chi_0}\qquad\text{ for any }\bm{c},
  \label{eq:def-Astar}
\end{equation}
for any \(\bm{c}\).
% with \(\bm{A}_*\) denoting the Piola transformation. 
Similarly, a dual
\(\operatorname{B}^*\) to the equilibrium operator \(\operatorname{B}\) in
\eqref{eq:strong} is defined with respect to \eqref{eq:conf-pair}:
\begin{equation}
  \left<\bm{\eta},\, \operatorname{B}_\chi\bm{\sigma}\right>_\chi = \left<\operatorname{B}_\chi^*\bm{\eta},\, \bm{\sigma}\right>_\chi.
  \label{eq:def-adjoint}
\end{equation}
% \tightlist
\item[A3]
The material strain measure \(\bm{e}\) is compatible with
\(\operatorname{B}^*_\chi\) through its directional derivative: 
\begin{equation}
  D\bm{e}\cdot\bm{u} \equiv \bm{A}\operatorname{B}^*_\chi \bm{u}.
  \label{eq:vstrain}
\end{equation}
%
\item[A4]
The constitutive response is defined between the material stress
\(\bm{s}\) and the material strain \(\bm{e}\) with material \(\bm{C}_s\) and
spatial \(\bm{C}_\sigma\) moduli defined by:
\begin{equation}
  D \bm{s}\cdot \bm{u} = \bm{C}_s \, D \bm{e}\cdot {\bm{u}}
  \qquad\text{ and }\qquad
  \bm{C}_\sigma \triangleq \bm{A}_*\bm{C}_s\bm{A}.
  \label{eq:def-moduli}
\end{equation}
\end{description}
\fi 
